\section{Method}
\label{sec:method}

\subsection{Pivot from the proposal}
\label{sec:pivot}

The original proposal assumed access to individual-level SPARK
genotypes, which would have allowed direct supervised training of a
gene-network deep learning model on per-individual case/control
labels. The two datasets that are publicly accessible without an IRB
review, however, are both \emph{summary statistics} (per-SNP effect
sizes and standard errors aggregated over all individuals): the
iPSYCH--PGC ASD November 2017 release and the Stein-Lab SPARK + iPSYCH
+ PGC meta-analysis. Individual-level SPARK genotypes are gated by
SFARI Base / dbGaP committee approval, which is incompatible with a
short coursework timeline.

We therefore reformulated the methodology as a \emph{sumstats-only}
pipeline that retains the spirit of the proposal -- improve PRS by
moving beyond independent additive scoring -- while operating purely
on aggregated effect-size estimates. The pivot has three concrete
consequences:

\begin{enumerate}[nosep,leftmargin=1.4em]
  \item We \emph{back out} an independent SPARK-only signal per SNP
        from the meta-analysis and discovery files (see below), and
        use that as the target for supervised learning.
  \item Both ML re-weighters (gradient-boosted and MLP) are trained at
        the SNP level rather than the individual level: each row is a
        SNP, the inputs are properties of that SNP, and the target is
        the derived SPARK-only score.
  \item The gene-network model pools SNP-level signal at the gene
        level over a STRING gene--gene interaction graph, then blends
        gene-level predictions back with each SNP's discovery signal.
\end{enumerate}

The four required deliverables of the proposal -- baseline PRS, data
preprocessing, model development, and evaluation against a baseline --
are all preserved.

\subsection{Data}

We use two publicly available ASD GWAS summary-statistics releases:
\begin{enumerate}[nosep,leftmargin=1.2em]
  \item \textbf{iPSYCH--PGC ASD Nov.~2017}~\cite{grove2019} with about
        46\,350 people (18\,381 cases, 27\,969 controls), reporting per
        SNP an rsID, GRCh37 chromosome and position, the two alleles,
        an imputation INFO score, the odds ratio, its standard error,
        and the two-sided $p$-value.
  \item \textbf{SPARK + iPSYCH + PGC meta-analysis}~\cite{matoba2020}
        with about 58\,794 people, reporting per SNP a chr:pos:alleles
        marker name, both alleles, the sample-size-weighted effect,
        its standard error, the $p$-value, the per-cohort direction
        string, and the total sample size.
\end{enumerate}
The two releases overlap by construction: the meta-analysis includes
the iPSYCH--PGC samples plus an additional SPARK contribution of about
12\,444 people.

\subsection{Quality control and harmonisation}

The two files are on different genome builds: the iPSYCH--PGC release
uses GRCh37 coordinates, whereas the Stein-Lab meta-analysis uses
GRCh38. A na\"{i}ve join on chromosome and position therefore matches
only about 73\,000 of the roughly 9\,million SNPs in either file. We
instead extract the rsID embedded in the meta-analysis
\texttt{MarkerName} field (the trailing \texttt{rs} number after the
final underscore) and join on rsID, keeping the iPSYCH GRCh37
coordinates as the canonical position. This recovers about
7.4\,million shared SNPs.

Both files are streamed with \texttt{polars}. Palindromic A/T and C/G
SNPs are removed because their alleles cannot be strand-aligned without
per-cohort allele frequencies. We further drop imputation \texttt{INFO}
below 0.6 and pairs with mismatched non-palindromic alleles, sign-flipping
the meta effect when alleles are swapped. After QC we retain
\textbf{6\,309\,737} SNPs (37\,101 in MHC). The major histocompatibility
complex (chr~6, 25--34\,Mb GRCh37) is flagged but kept; downstream
training scripts exclude it from the train/validation folds.

\subsection{Independent SPARK-only target}
\label{sec:sparkz}

The two summary-statistics files are not statistically independent
because the meta-analysis includes the iPSYCH samples; using one as
training and the other as test would therefore leak information. To
obtain an independent target without individual-level genotypes, we use
the standard relationship between a combined meta-analysis Z-score and
the contributing cohorts. Intuitively, the meta Z is a
sample-size-weighted blend of the iPSYCH discovery Z and the
additional SPARK contribution. If we know the total sample sizes for
the meta-analysis, the iPSYCH-only analysis, and the extra SPARK
slice, we can \emph{solve} for the SPARK-only Z at each SNP by
subtracting the appropriately scaled iPSYCH part from the scaled meta
part and re-scaling by the SPARK sample size. In our setting the
iPSYCH-only analysis has 46\,350 people, the meta-analysis 58\,794, and
the implied SPARK-only piece 12\,444.

We verify this step by \emph{round-tripping}: plugging the recovered
SPARK Z back into the same blending rule must reproduce the published
meta Z at every SNP. In practice the maximum reconstruction error is
below $10^{-15}$, i.e.\ numerical noise only.

\subsection{Baseline 1: Pruning and Thresholding (P+T)}

For each of several $p$-value cut-offs (from genome-wide significance
down to ``keep everything''), we LD-clump the discovery sumstats with
PLINK~1.9 (\texttt{--clump} with correlation under 0.1 in a 250\,kb
window) using the 1000 Genomes phase~3 EUR reference panel. SNPs that
survive clumping keep their discovery log-odds as the PRS weight; all
other SNPs get weight zero. When PLINK or the LD reference are
unavailable, the implementation falls back to a greedy distance-based
prune so the pipeline still runs end-to-end.

\subsection{Baseline 2: Empirical-Bayes shrinkage (PRS-CS-style)}

PRS-CS~\cite{ge2019prscs} is a Bayesian method that shrinks SNP effects
using linkage-disequilibrium structure from a reference panel. We did
not download the full LD reference; instead we use a simple
closed-form shrinkage that pulls each discovery effect toward zero when
its Z-score is small, and leaves larger Z-scores more intact. This
stands in for PRS-CS in our environment and preserves the same
qualitative comparison against the ML methods.

\subsection{Feature-based ML re-weighter}
\label{sec:mlrew}

For every harmonised SNP we build a fixed list of \textbf{14 input
features}: the discovery Z-score and its absolute value, the imputation
INFO score, and summaries (mean, max, count) of absolute Z in a
$\pm 500$\,kb window around the SNP (a simple stand-in for local LD
context). When annotation files are present, the same vector also
includes distance to the nearest gene, SFARI flags, STRING degree,
GTEx brain eQTL strength, ENCODE chromatin state, and conservation; in
our run several of these were missing and were filled with zeros.

The training target at each SNP is the SPARK-only Z from the previous
subsection. We train two models with a chromosome-based split
(Section~\ref{sec:experiments}): \textbf{XGBoost} (gradient boosting,
depth and learning rate set for stable validation error) on the full
training chromosomes, and a \textbf{small neural network} (two hidden
layers of 128 and 64 units) via scikit-learn's \texttt{MLPRegressor},
trained on a random 250\,k-row subset of training SNPs for speed. To
turn predicted Z back into an effect size we multiply by the discovery
standard error, as usual for Z-based regression.

\subsection{Gene-network GNN re-weighter}
\label{sec:gnn}

We map each SNP to its nearest gene (within $\pm 10$\,kb of the gene
body) and average SNP-level features within each gene. A gene graph
comes from STRING~v12 high-confidence edges among genes that have at
least one SNP. A two-layer graph neural network predicts the average
SPARK-only Z per gene, then we blend each gene's prediction with that
SNP's own discovery Z (equal 50--50 weight in the default setup). When
the full graph library is unavailable we use a lighter Laplacian
smoothing fallback on the same graph. In our run, missing gene
annotation files meant no SNPs mapped to genes, so this path
collapsed to the discovery baseline for reporting purposes.
